%%%PAGE 242 rot=0 legibility=good summary=六道题：1.两绳夹角不变的动态平衡 2.平衡到向右匀加速绳张力 3.斜面平抛速度比较 4.转板带电粒子 5.双球相似三角形 6.靠墙杆的受力（力矩法与分解法）
%%%BLOCK id=242-01 ch=02 sec=动态平衡·两绳夹角不变 kind=problem alt=none cont=none y=0.00-0.24
\begin{problem}[1]
%FIG p242_a 0.07 0.00 0.235 0.072
\notefig[0.25]{figures/p242_a.png}
%FIG p242_b 0.30 0.00 0.40 0.05
\notefig[0.25]{figures/p242_b.png}
%FIG p242_c 0.12 0.075 0.29 0.15
\notefig[0.25]{figures/p242_c.png}
\begin{solution}
\[ \begin{cases} F_1\sin\alpha=F_2\sin(\theta-\alpha)\\ F_1\cos\alpha+F_2\cos(\theta-\alpha)=G \end{cases} \]
\[ F_1=G\,\frac{\sin(\theta-\alpha)}{\sin\theta}\qquad F_2=G\,\frac{\sin\alpha}{\sin\theta}\qquad\theta\text{不变} \]
\begin{itemize}
\item $\alpha$从$90^\circ\to0^\circ$
\item $\theta-\alpha$\unclear{在}\ $120^\circ$\ 从$30^\circ\sim120^\circ$ % CHECK: 此处三行斜排：“θ−α ?”“120°”“从30°~120°”，120°应指θ=120°
\item $F_1$先$\uparrow$后$\downarrow$
\item $F_2\downarrow\to0$
\end{itemize}
\end{solution}
\end{problem}
%%%END
%%%BLOCK id=242-02 ch=03 sec=瞬时与加速·绳张力变化 kind=problem alt=02 cont=none y=0.20-0.36
\begin{problem}[2]
%FIG p242_d 0.095 0.205 0.24 0.30
\notefig[0.25]{figures/p242_d.png}
由平衡到向右匀加速
\begin{solution}
原
$\begin{cases} mg=T_a\cos\theta\\ T_b=T_a\sin\theta \end{cases}$
\[ T_a\sin\theta-T_b=ma\qquad\theta\text{不变}\quad T_a\text{不变}\quad\therefore T_b\downarrow \]
\end{solution}
\end{problem}
%%%END
%%%BLOCK id=242-03 ch=04 sec=平抛·斜面上的平抛 kind=problem alt=none cont=none y=0.36-0.56
\begin{problem}[3、平抛]
%FIG p242_e 0.075 0.386 0.375 0.507
\notefig[0.3]{figures/p242_e.png}
速度反向延长线\\交于位移中点
\begin{solution}
\begin{align*}
&t_A<t_B<t_C\\
&x_A=x_B=x_C\\
&\therefore V_{A0}>V_{B0}>V_{C0}
\end{align*}
\[ h=\frac{x}{2}\,\frac{1}{\tan\theta}\qquad t=\sqrt{\frac{2h}{g}}\qquad V_y=\sqrt{2gh}\qquad V_x=\tan\theta\sqrt{2gh} \]
\[ \therefore V=\sqrt{V_x^2+V_y^2}=\sqrt{2gh\,(1+\tan^2\theta)}=\sqrt{xg\Bigl(\frac{1}{\tan\theta}+\frac{1}{\tan\theta}\Bigr)} \] % CHECK: 按推导括号内第二项应为 tanθ，原稿写作 1/tanθ
代入 $60^\circ$ $45^\circ$ $30^\circ$ 得 $V_A=V_C>V_B$
\end{solution}
\end{problem}
%%%END
%%%BLOCK id=242-04 ch=09 sec=带电粒子·平行板间平衡(转板) kind=problem alt=03 cont=none y=0.565-0.69
\begin{problem}[4]
%FIG p242_f 0.115 0.562 0.40 0.69
\notefig[0.3]{figures/p242_f.png}
转板
\begin{solution}
原\quad $mg=qE=q\dfrac{U}{d}$

现\quad $\dfrac{U}{d}q=\dfrac{U}{d\cos\theta}=\dfrac{mg}{\cos\theta}$ % CHECK: 第二项分子无 q（原稿如此）

合力水平向左
\end{solution}
\end{problem}
%%%END
%%%BLOCK id=242-05 ch=02 sec=共点力平衡·相似三角形 kind=problem alt=none cont=none y=0.70-0.865
\begin{problem}[5]
%FIG p242_g 0.07 0.70 0.225 0.865
\notefig[0.25]{figures/p242_g.png}
\begin{solution}
\[ \frac{T}{l_1}=\frac{m_1g}{h}\qquad\frac{T}{l_2}=\frac{m_2g}{h} \]
\[ \therefore m_1l_1=m_2l_2 \]
\end{solution}
\end{problem}
%%%END
%%%BLOCK id=242-06 ch=02 sec=共点力平衡·靠墙杆的受力 kind=problem alt=none cont=none y=0.865-1.00
\begin{problem}[6]
%FIG p242_h 0.11 0.865 0.245 0.985
\notefig[0.25]{figures/p242_h.png}
\begin{solution}
法1\ 用杆\unclear{力矩}平衡 % CHECK: “力矩”二字笔迹潦草，仅辨认出大意
\[ \underset{\uparrow}{F_1}\,OA=G\,O\underset{\uparrow}{B} \] % 原稿 F_1 与 B 的下方各有一个向上的小箭头
\[ N_2=G\qquad f=F_1 \]
地对木杆力\unclear{即}$\uparrow$ $\vec N+\vec f$

%FIG p242_i 0.46 0.88 0.585 0.96
\notefig[0.25]{figures/p242_i.png}

法2\ \unclear{竖水平分} % CHECK: 末字不清（也可能是“衡”），大意应为竖直/水平方向分解
%FIG p242_j 0.72 0.915 0.885 0.98
\notefig[0.25]{figures/p242_j.png}
\end{solution}
\end{problem}
%%%END
%%%PAGEEND 242
